\section{Scales and conventions}\label{sec:setup}

We use the Hamiltonian of Section~\ref{sec:introduction} also on labelled
coordinates, and suppress $L$ when it is fixed. We choose an everywhere
nonnegative bounded Borel radial representative of $v$ vanishing outside
its range ball. Changing this representative on a null set changes neither
the operator nor the path integrals: at interior times the free relative
positions have densities. Write $R_v$ for a range bound, and use
$\|v\|_1,\|v\|_\infty$ on $\R^3$. In formulas on the torus,
$v$ abbreviates its periodization $v_L$; infinite-space norms always
refer to the original potential.

The orthogonal symmetrizer is included exactly once in a trace kernel.
If $k_t$ is the labelled $N$-particle kernel, then
\[
 k_t^{\mathrm{sym}}(X,Y)=\frac1{N!}\sum_{\pi\in\mathfrak S_N}
                         k_t(X,\pi Y),\qquad
 Z_N=\int k_\beta^{\mathrm{sym}}(X,X)\dd X.
\]
Thus integration over labelled coordinates carries no further factorial.
We work with $\rho/2\le N/V\le2\rho$, $0<T\le\rho^2$, and
$\beta=T^{-1}$. Every volume restriction below is imposed at fixed $\rho$,
uniformly over this temperature interval. Constants $C,c>0$ may change
between occurrences and depend on $v$ and previously fixed cutoffs,
but not on $\rho,L,N,T$. The final order of those cutoff choices is
recorded in Section~\ref{sec:conclusion}.

We use a periodic grid \(\mathcal G=(\mathbb Z/n\mathbb Z)^3\) of equal cubes (fix the assignment at faces), with
\[
 n=\left\lfloor\frac{L}{K\rho^{-1/2}}\right\rfloor,\qquad
 s=L/n\in[K\rho^{-1/2},2K\rho^{-1/2}],\qquad
 M=n^3,\quad m=N/M .
\]
Here \(K\ge1\) is a fixed large constant chosen below. Thus \(\rho s^3/2\le m\le 2\rho s^3\). 

A union \(U\) of \(u\ge1\) cells has volume \(|U|=u s^3\). A neighborhood out to distance \(h s\), \(h\ge0\), can be covered by at most \(C(h+2)^3 u\) cells, including on the torus. Physical distances \(d(\cdot,\cdot)\) are Euclidean torus distances unless indicated otherwise. Cell index distances will use the periodic sup norm. For a function \(p\) of position put \(n_p(X)=\sum_i p(x_i)\) on configurations (also the multiplication operator); similarly \(n_U\) for counting particles in \(U\).

